% ============================================================
% Temario — ruta de sesiones por grado y periodo
% Generado desde content/guias/*.yaml (180 guias)
% ============================================================

\needspace{6\baselineskip}
\bigskip
\noindent{\large\color{vinoInst}\textbf{Grado Sexto (6°)}}
\vspace{1pt}\par\noindent{\color{vinoInst}\rule{\textwidth}{1pt}}
\vspace{3pt}

\needspace{5\baselineskip}
\noindent\textbf{Periodo 1 · Comunicación e identidad digital}\hfill{\footnotesize\color{vinoInst}10 guías}\par
\begin{enumerate}[leftmargin=2.6em, itemsep=0.6pt, topsep=2pt, parsep=0pt, label=\footnotesize\arabic*.]
\footnotesize
  \item Bienvenida a Tecnología --- cómo aprenderemos este año
  \item Mi identidad digital --- cómo me presento en la red sin perderme
  \item Mi correo institucional --- la primera dirección que firmas como tuya
  \item Netiqueta --- los 10 acuerdos para llevarte bien en internet
  \item Historia de los medios --- del pregonero al chat de WhatsApp
  \item Mi huella digital --- el rastro que dejas sin darte cuenta
  \item Privacidad y datos personales --- cerrar las puertas con criterio
  \item Riesgos digitales --- reconocer la trampa antes de caer
  \item Comunicación responsable --- los 3 filtros antes de publicar
  \item Mi presentación digital --- 2 minutos para mostrar lo que aprendí
\end{enumerate}

\needspace{5\baselineskip}
\noindent\textbf{Periodo 2 · Hardware, software y la máquina por dentro}\hfill{\footnotesize\color{vinoInst}10 guías}\par
\begin{enumerate}[leftmargin=2.6em, itemsep=0.6pt, topsep=2pt, parsep=0pt, label=\footnotesize\arabic*.]
\footnotesize
  \item Apertura periodo 2 --- las máquinas que nos rodean por dentro
  \item ¿Qué es un computador? --- las 4 funciones básicas de toda máquina inteligente
  \item El gabinete por dentro --- conocer las piezas como las conoce el relojero
  \item Periféricos --- las manos, los ojos y la voz del computador
  \item Hardware vs software --- el cuerpo y las instrucciones
  \item El sistema operativo --- el gerente que coordina toda la máquina
  \item Archivos y carpetas --- ordenar mi biblioteca digital
  \item Mantenimiento básico --- cuidar el computador como cuida el relojero
  \item Postura y vida útil --- cuidarte tú mientras usas la máquina
  \item Mi computador ideal --- diseñar la máquina que yo necesito
\end{enumerate}

\needspace{5\baselineskip}
\noindent\textbf{Periodo 3 · Escribir y buscar con criterio}\hfill{\footnotesize\color{vinoInst}10 guías}\par
\begin{enumerate}[leftmargin=2.6em, itemsep=0.6pt, topsep=2pt, parsep=0pt, label=\footnotesize\arabic*.]
\footnotesize
  \item Apertura periodo 3 --- el oficio de escribir y buscar con criterio
  \item ¿Qué es un procesador de texto? --- mi primera página escrita en digital
  \item Formato básico --- darle voz visual a lo que escribo
  \item Listas, viñetas y numeración --- ordenar la información para que se entienda
  \item Tablas e imágenes --- mostrar datos y agregar voz visual
  \item Encabezado, pie de página y números --- terminar el documento como un profesional
  \item ¿Qué es internet? --- cómo viajan los datos para llegar a tu pantalla
  \item Buscar bien en internet --- encontrar lo que necesito en segundos
  \item Evaluar fuentes --- distinguir lo verdadero de lo inventado en internet
  \item Mi primer documento informativo --- producto que reúne todo lo aprendido
\end{enumerate}

\needspace{6\baselineskip}
\bigskip
\noindent{\large\color{vinoInst}\textbf{Grado Séptimo (7°)}}
\vspace{1pt}\par\noindent{\color{vinoInst}\rule{\textwidth}{1pt}}
\vspace{3pt}

\needspace{5\baselineskip}
\noindent\textbf{Periodo 1 · Trabajo colaborativo en la nube (Microsoft 365)}\hfill{\footnotesize\color{vinoInst}10 guías}\par
\begin{enumerate}[leftmargin=2.6em, itemsep=0.6pt, topsep=2pt, parsep=0pt, label=\footnotesize\arabic*.]
\footnotesize
  \item Apertura periodo 1 --- vamos a hacer minga digital
  \item Microsoft 365 --- conocer las herramientas del taller digital
  \item OneDrive --- mi nube personal y mi cuarto de archivos
  \item Word, Excel y PowerPoint en línea --- las 3 herramientas estrella
  \item Compartir y permisos --- abrir mi taller a invitados con criterio
  \item Coautoría asincrónica --- 3 personas escribiendo al mismo tiempo
  \item Comentarios y sugerencias --- revisar el trabajo ajeno con respeto
  \item Versiones e historial --- el sistema recuerda todo lo que escribiste
  \item Outlook y Teams --- comunicarme con voz profesional
  \item Cosecha del periodo --- mi primer proyecto colaborativo
\end{enumerate}

\needspace{5\baselineskip}
\noindent\textbf{Periodo 2 · Algoritmia y pensamiento computacional (Scratch)}\hfill{\footnotesize\color{vinoInst}10 guías}\par
\begin{enumerate}[leftmargin=2.6em, itemsep=0.6pt, topsep=2pt, parsep=0pt, label=\footnotesize\arabic*.]
\footnotesize
  \item Apertura periodo 2 --- pensar como teje la abuela
  \item ¿Qué es un algoritmo? --- la receta paso a paso
  \item Pensamiento computacional --- las 4 técnicas del oficio
  \item Diagramas de flujo --- dibujar el algoritmo
  \item Variables --- cajas con etiqueta que guardan información
  \item Condicionales --- el si... entonces de los algoritmos
  \item Bucles --- repeticiones inteligentes en los algoritmos
  \item Scratch --- mi primer programa con bloques
  \item Scratch · narrativa interactiva --- el cuento que el lector escoge
  \item Cosecha del periodo 2 --- mi primer programa Scratch completo
\end{enumerate}

\needspace{5\baselineskip}
\noindent\textbf{Periodo 3 · Inteligencia artificial con criterio}\hfill{\footnotesize\color{vinoInst}10 guías}\par
\begin{enumerate}[leftmargin=2.6em, itemsep=0.6pt, topsep=2pt, parsep=0pt, label=\footnotesize\arabic*.]
\footnotesize
  \item Apertura periodo 3 --- la IA como el nuevo consejero del barrio
  \item ¿Qué es la IA? --- definición, historia y casos cotidianos
  \item Tipos de IA --- estrecha, general, modelos especializados
  \item Cómo aprende una IA --- datos, entrenamiento y el problema del sesgo
  \item Modelos de lenguaje (LLMs) --- ChatGPT, Claude, Gemini
  \item Prompting básico --- cómo hablarle a la IA para que te responda bien
  \item Prompting avanzado --- técnicas que diferencian al usuario experto
  \item Ética de la IA --- sesgo, privacidad, derechos y deepfakes
  \item IA y mis tareas --- política personal de uso responsable en estudios
  \item Cosecha del periodo 3 --- mi proyecto con IA responsable
\end{enumerate}

\needspace{6\baselineskip}
\bigskip
\noindent{\large\color{vinoInst}\textbf{Grado Octavo (8°)}}
\vspace{1pt}\par\noindent{\color{vinoInst}\rule{\textwidth}{1pt}}
\vspace{3pt}

\needspace{5\baselineskip}
\noindent\textbf{Periodo 1 · Lógica computacional avanzada}\hfill{\footnotesize\color{vinoInst}10 guías}\par
\begin{enumerate}[leftmargin=2.6em, itemsep=0.6pt, topsep=2pt, parsep=0pt, label=\footnotesize\arabic*.]
\footnotesize
  \item Phronesis con datos --- preguntar antes de calcular
  \item Tipos de datos y formato de celdas en Excel
  \item Tablas de datos --- campos, registros y limpieza con bitácora
  \item SUMA, PROMEDIO, MAX, MIN --- 4 funciones que dirigen decisiones
  \item Referencias relativas y absolutas en Excel
  \item Fórmulas compuestas --- operadores, porcentajes y jerarquía PEMDAS
  \item Gráficos básicos --- barras, líneas y circulares con honestidad visual
  \item Formato condicional y validación de datos --- alertas visibles
  \item Mini-estudio de datos del entorno escolar
  \item Sustentación del mini-estudio --- 4 minutos de phronesis comunicativa
\end{enumerate}

\needspace{5\baselineskip}
\noindent\textbf{Periodo 2 · Datos como tablas: pensar con estructura}\hfill{\footnotesize\color{vinoInst}10 guías}\par
\begin{enumerate}[leftmargin=2.6em, itemsep=0.6pt, topsep=2pt, parsep=0pt, label=\footnotesize\arabic*.]
\footnotesize
  \item Lógica avanzada --- if-else con AND, OR y NOT
  \item Tablas de verdad --- AND, OR, NOT como reglas universales
  \item Algoritmos --- del pseudocódigo al diagrama de flujo
  \item Sensores en MakeCode --- el micro:bit que siente
  \item Actuadores --- LEDs, sonido y coreografía
  \item Sensores con variables, umbrales y calibración
  \item Depuración --- buscar el error como el relojero busca el tic-tac
  \item Alertas con lógica compuesta --- validación con escenarios
  \item Proyecto de monitoreo --- el micro:bit como vigía
  \item Sustentación del proyecto técnico --- honestidad del oficio
\end{enumerate}

\needspace{5\baselineskip}
\noindent\textbf{Periodo 3 · Diseño visual y comunicación digital}\hfill{\footnotesize\color{vinoInst}10 guías}\par
\begin{enumerate}[leftmargin=2.6em, itemsep=0.6pt, topsep=2pt, parsep=0pt, label=\footnotesize\arabic*.]
\footnotesize
  \item Multimedia --- diseño visual y jerarquía de información
  \item Animaciones e hipervínculos --- narrativa interactiva
  \item Narrativa de alto impacto --- estructura TED-style
  \item Edición de imagen --- ética visual y derechos de imagen
  \item Audio --- captura, edición y formatos
  \item Ciberbullying --- prevención y Línea 141
  \item Grooming y sexting --- riesgos digitales y Ley 1336
  \item Estética de la liberación --- cultura popular y producción ética
  \item Proyecto integrador comunitario --- voz para el cambio social
  \item Sustentación pública y portafolio digital del grado 8
\end{enumerate}

\needspace{6\baselineskip}
\bigskip
\noindent{\large\color{vinoInst}\textbf{Grado Noveno (9°)}}
\vspace{1pt}\par\noindent{\color{vinoInst}\rule{\textwidth}{1pt}}
\vspace{3pt}

\needspace{5\baselineskip}
\noindent\textbf{Periodo 1 · Técnica, ciencia y tecnología}\hfill{\footnotesize\color{vinoInst}10 guías}\par
\begin{enumerate}[leftmargin=2.6em, itemsep=0.6pt, topsep=2pt, parsep=0pt, label=\footnotesize\arabic*.]
\footnotesize
  \item ¿Qué es técnica? --- del cuerpo al instrumento
  \item La rueda y el agua --- máquinas simples del campo colombiano
  \item La balanza y el peso --- medir como acto político
  \item Línea del tiempo de la técnica --- del fuego a la imprenta
  \item Revolución industrial --- máquinas que cambiaron el trabajo
  \item La imprenta --- escribir muchas veces lo mismo
  \item La electricidad --- luz, motor, comunicación
  \item Era digital --- del ábaco al chip
  \item Tecnologías propias --- saberes que el mundo aprende ahora
  \item Cosecha P1 --- manifiesto del técnico crítico
\end{enumerate}

\needspace{5\baselineskip}
\noindent\textbf{Periodo 2 · Diseño editorial y maquetación digital}\hfill{\footnotesize\color{vinoInst}10 guías}\par
\begin{enumerate}[leftmargin=2.6em, itemsep=0.6pt, topsep=2pt, parsep=0pt, label=\footnotesize\arabic*.]
\footnotesize
  \item ¿Qué es diseñar? --- orden, jerarquía, intención
  \item La página y la cuadrícula --- primer principio del orden
  \item Tipografía --- la voz visible de las palabras
  \item Color en editorial --- contraste, jerarquía, accesibilidad
  \item Imagen y texto --- el ritmo de la lectura
  \item Herramientas digitales --- Figma, Canva o Adobe Express
  \item Accesibilidad lectora --- quién NO puede leer mi revista y por qué
  \item Edición y corrección --- el ojo crítico
  \item Maquetar la revista completa --- flujo de trabajo
  \item Cosecha P2 --- sustentación pública de la revista
\end{enumerate}

\needspace{5\baselineskip}
\noindent\textbf{Periodo 3 · Datos, estadística y visualización}\hfill{\footnotesize\color{vinoInst}10 guías}\par
\begin{enumerate}[leftmargin=2.6em, itemsep=0.6pt, topsep=2pt, parsep=0pt, label=\footnotesize\arabic*.]
\footnotesize
  \item ¿Qué es un dato? --- recolección con propósito
  \item Tablas y registros --- la columna como tipo, la fila como caso
  \item Fórmulas básicas --- SUMA, PROMEDIO, MAX, MIN, CONTAR
  \item Filtros y orden --- preguntar a los datos
  \item Tablas dinámicas --- agrupar y comparar
  \item Gráficos --- la imagen de los datos
  \item Limpieza de datos --- el oficio del archivero
  \item Visualización honesta --- gráficos que mienten
  \item Insights --- del dato a la decisión
  \item Cosecha P3 --- sustentación del mini-estudio
\end{enumerate}

\needspace{6\baselineskip}
\bigskip
\noindent{\large\color{vinoInst}\textbf{Grado Décimo (10°)}}
\vspace{1pt}\par\noindent{\color{vinoInst}\rule{\textwidth}{1pt}}
\vspace{3pt}

\needspace{5\baselineskip}
\noindent\textbf{Periodo 1 · Edición digital y autoría con IA responsable}\hfill{\footnotesize\color{vinoInst}10 guías}\par
\begin{enumerate}[leftmargin=2.6em, itemsep=0.6pt, topsep=2pt, parsep=0pt, label=\footnotesize\arabic*.]
\footnotesize
  \item El oficio del editor --- qué hace y qué no hace
  \item Concebir tu libro --- género, tema, audiencia y escaleta
  \item Prompting técnico para escritura --- las 5 partes del prompt profesional
  \item Iterar y refinar texto con IA --- del borrador al capítulo terminado
  \item Estructura editorial --- arco narrativo, ritmo y capítulos
  \item Derechos de autor y propiedad intelectual con IA --- Ley 23 de 1982 y obras derivadas
  \item Portada con IA generativa de imagen --- herramientas gratuitas
  \item Diagramación con herramientas gratuitas --- Canva, Google Docs, LibreOffice
  \item Producción final del libro --- 80 páginas compiladas, revisadas y firmadas
  \item Sustentación + feria del libro escolar
\end{enumerate}

\needspace{5\baselineskip}
\noindent\textbf{Periodo 2 · Comunicación profesional con IA}\hfill{\footnotesize\color{vinoInst}10 guías}\par
\begin{enumerate}[leftmargin=2.6em, itemsep=0.6pt, topsep=2pt, parsep=0pt, label=\footnotesize\arabic*.]
\footnotesize
  \item 4 tipos de texto profesional --- informe, comercial, técnico, estudio de mercado
  \item Estructura del informe técnico --- introducción, desarrollo, conclusiones
  \item Google Docs para informes profesionales --- estilos, índice y citas
  \item Markdown --- el formato del oficio digital portable
  \item Encuestas con phronesis --- Google Forms más IA para preguntar bien
  \item Análisis cualitativo con IA --- qué hace y qué no
  \item Informe comercial --- para audiencia real
  \item Correo profesional con IA --- asistente, no reemplazo
  \item Mini-proyecto --- informe técnico con datos reales y IA documentada
  \item Sustentación del periodo 2 --- defensa del oficio
\end{enumerate}

\needspace{5\baselineskip}
\noindent\textbf{Periodo 3 · Microemprendimiento y gestión digital}\hfill{\footnotesize\color{vinoInst}10 guías}\par
\begin{enumerate}[leftmargin=2.6em, itemsep=0.6pt, topsep=2pt, parsep=0pt, label=\footnotesize\arabic*.]
\footnotesize
  \item Ofimática con IA --- el copiloto del oficio digital cotidiano
  \item Contabilidad básica --- ingresos, egresos y balance personal o familiar
  \item Google Sheets + IA para contabilidad personal o microempresa
  \item Flujo de caja y punto de equilibrio --- proyectar con IA
  \item Visualización de datos contables con IA --- gráficos honestos
  \item Encuestas y análisis de mercado con IA
  \item Plan de negocio con IA --- Business Model Canvas asistido
  \item Comunicación empresarial profesional con IA
  \item Proyecto integrador --- micro-emprendimiento con plan completo
  \item Sustentación final + feria de microempresas escolares
\end{enumerate}

\needspace{6\baselineskip}
\bigskip
\noindent{\large\color{vinoInst}\textbf{Grado Undécimo (11°)}}
\vspace{1pt}\par\noindent{\color{vinoInst}\rule{\textwidth}{1pt}}
\vspace{3pt}

\needspace{5\baselineskip}
\noindent\textbf{Periodo 1 · Presencia digital empresarial con IA}\hfill{\footnotesize\color{vinoInst}10 guías}\par
\begin{enumerate}[leftmargin=2.6em, itemsep=0.6pt, topsep=2pt, parsep=0pt, label=\footnotesize\arabic*.]
\footnotesize
  \item Mi marca digital --- del oficio heredado a la huella propia
  \item Identidad visual --- del grafismo ancestral a la paleta propia
  \item Sitio web personal --- fundamentos de la casa digital
  \item Plantillas web --- elegir y adaptar sin perderse uno mismo
  \item Copywriting con IA --- escribir con criterio, no a ciegas
  \item SEO básico --- hacerme visible sin engaños
  \item Imagen propia --- cómo se ve mi marca en foto y video
  \item Perfiles profesionales --- LinkedIn y portafolio coherentes
  \item Métricas básicas --- quién entra a mi sitio y qué busca
  \item Lanzamiento P1 --- mi presencia digital sale al mundo
\end{enumerate}

\needspace{5\baselineskip}
\noindent\textbf{Periodo 2 · Automatización y gestión de procesos}\hfill{\footnotesize\color{vinoInst}10 guías}\par
\begin{enumerate}[leftmargin=2.6em, itemsep=0.6pt, topsep=2pt, parsep=0pt, label=\footnotesize\arabic*.]
\footnotesize
  \item Mapeo de procesos --- ver dónde se pierde el tiempo
  \item Diagramas de flujo --- lenguaje BPMN básico
  \item Formularios digitales --- captura de información sin papel
  \item Hojas de cálculo como base de datos
  \item Automatización con triggers --- Zapier, Make, n8n
  \item IA en mensajería --- clasificar y responder con criterio
  \item Chatbots sin código --- flujos de conversación
  \item KPIs --- métricas que importan en un proceso
  \item Mejora continua --- Kaizen aplicado a procesos digitales
  \item Caso real --- automatizar un proceso del colegio o la casa
\end{enumerate}

\needspace{5\baselineskip}
\noindent\textbf{Periodo 3 · Microempresa con MVP, datos y ética}\hfill{\footnotesize\color{vinoInst}10 guías}\par
\begin{enumerate}[leftmargin=2.6em, itemsep=0.6pt, topsep=2pt, parsep=0pt, label=\footnotesize\arabic*.]
\footnotesize
  \item Ideación --- un problema real de mi comunidad
  \item Validación --- hablar antes de construir
  \item Modelo de negocio --- el lienzo de tu proyecto
  \item MVP digital --- una versión imperfecta para probar
  \item Datos del MVP --- recoger evidencia con rigor
  \item Pitch --- contar la historia de tu proyecto
  \item Presupuesto --- cuánto siembras, cuánto cosechas
  \item Ética y legal --- la palabra empeñada en lo digital
  \item Versión final del MVP
  \item Sustentación pública --- el aprendiz frente a los mayores
\end{enumerate}

